\section{判据边界的正例与反例}
\label{app:examples-counterexamples}

本附录给出五类判据边界：循环链的微观尺度障碍、无群等变性的精确闭合、
逐环境宏观核与共享宏观核的量词差异，以及闭合证书事件与群作用分离事件
两个方向的反例。

\subsection{循环链的微观尺度障碍}

\begin{counterexample}[循环链]
\label{counter:cycle}
对 $N\ge2$，令 $B_N=\mathbb Z/N\mathbb Z$，确定性 kernel 为
$K_N(x,x+1)=1$，任务为 $g_N(x)=\one_{\{x=0\}}$。若 $0\le\tau<1/2$，
任何任务保真且闭合缺陷不超过 $\tau$ 的分区都只能是离散分区。因此
\[
c_{K_N}^{\tau}(g_N)=N,
\]
因此该系统族的相对状态复杂度保持为 $1$，不满足 LC。
\end{counterexample}

\begin{proof}
若 $s(x)=s(y)$，两条投影后的一步分布 $\delta_{s(x+1)}$ 与
$\delta_{s(y+1)}$ 都在 TV 距离 $\tau$ 内接近同一宏观行。因此它们之间的
距离不超过 $2\tau<1$。两个 Dirac 分布的 TV 距离只能是 $0$ 或 $1$，故
$s(x+1)=s(y+1)$；迭代得到 $s(x+t)=s(y+t)$ 对所有 $t$ 成立。若
$x\ne y$，可选 $t$ 使其中一个状态等于 $0$、另一个不等于 $0$，与任务
保真矛盾。
\end{proof}

\subsection{无群等变性的精确闭合}

\begin{example}[无群等变性的闭合商]
\label{ex:closure-without-symmetry}
取 $B=\{1,2,3\}$，分区 $\pi=\{\{1,2\},\{3\}\}$，并令
\[
K=\begin{pmatrix}
0.1&0.6&0.3\\
0.2&0.5&0.3\\
0.1&0.2&0.7
\end{pmatrix}.
\]
前两行到两个分区块的概率都为 $(0.7,0.3)$，所以 $\pi$ 精确闭合。交换
状态 $1,2$ 的完整 kernel 行不同。这说明群等变性提供一种充分构造，精确闭合
也可以由其他结构产生。
\end{example}

\subsection{逐环境宏观核与共享宏观核的量词差异}

\begin{counterexample}[两个量词次序]
\label{counter:environment-specific-versus-shared}
令 $B=\{0,1\}\times\{a,b\}$，$Y=\{0,1\}$，$q(y,z)=y$，任务取
$g=q$。构造两个 kernels：$K^{(0)}$ 的下一宏观状态恒为 $0$，
$K^{(1)}$ 的下一宏观状态恒为 $1$；纤维内第二坐标可按任意固定规则更新。
则 $q$ 对两者都精确闭合，分别对应
$L^{(0)}(y,\cdot)=\delta_0$ 与 $L^{(1)}(y,\cdot)=\delta_1$，故
\[
\sup_{i\in\{0,1\}}\min_L
d_\infty(K^{(i)}Q_q,Q_qL)=0.
\]
任意共享 $L$ 的每一行都必须同时逼近 $\delta_0$ 与 $\delta_1$。由三角
不等式其最坏误差至少为 $1/2$，而取均匀分布行可以达到该值。因此
\[
\min_L\sup_{i\in\{0,1\}}
d_\infty(K^{(i)}Q_q,Q_qL)=\frac12.
\]
\end{counterexample}

\subsection{闭合证书成立而群作用分离失败}

\begin{counterexample}[精确轨道压缩与单一群不变分布]
\label{counter:certificate-without-separation}
取第~\ref{sec:fault-repair}节中 $a,b\in(0,1)$ 的故障--修复链，并令
$G_N=H_N=\mathfrak S_N$ 按坐标置换作用。把所有对象放在一点评率空间上。
任务在轨道上不变，动力学精确 $G_N$-等变，而且轨道就是故障数水平集，故
\[
\eps_{H_N}(\{K_{N,a,b}\})=0,
\qquad
\frac{|B_N/H_N|}{|B_N|}=\frac{N+1}{2^N}\longrightarrow0.
\]
所以附录~\ref{app:structural-certificates}的闭合证书事件可取
$a_N=(N+1)/2^N$、$\eps_N=0$ 而恒成立。

$H_N=G_N$ 不满足群作用分离事件要求的 $H_N\subsetneq G_N$。更实质地，
$a,b>0$ 时有限链不可约，因而有唯一平稳分布；精确 $G_N$-等变性迫使该分布
在整个 $G_N$ 下不变。因此在
$\beta_N=\gamma_N=\kappa_N=0$、$\ell_N>0$ 的精确阈值下，群外分离条件
失败。该系统同时具有精确闭合轨道商和单一群不变平稳分布。
\end{counterexample}

\subsection{群作用分离成立而闭合压缩证书失败}

\begin{counterexample}[精确群作用分离与恒等任务]
\label{counter:separation-without-certificate}
令
\[
B_N=\{+,-\}\times\{1,\ldots,N\},
\qquad
g_N=\Id_{B_N},
\]
$G_N=\mathbb Z_2$ 交换第一坐标，$H_N=\{e\}$，
$K_N=\Id_{B_N}$，$\mu_N=\delta_{(+,1)}$，并仍取一点评率空间。此时
$K_N$ 对 $G_N$ 精确等变，$\mu_N$ 精确平稳，且
\[
\kappa_{H_N}(\mu_N)=0,
\qquad
\lambda_{H_N}(\mu_N)=1.
\]
因此群作用分离事件可取
$\beta_N=\gamma_N=\kappa_N=0$、$\ell_N=1$ 而恒成立。可是恒等任务迫使
每个任务保真表示都是单射，尤其
\[
c_{\{K_N\}}^\tau(g_N)=|B_N|
\qquad\text{对每个 }\tau\ge0,
\]
且 $|B_N/H_N|/|B_N|=1$。故任取 $a_N<1$，闭合证书事件都失败。该例表明
群作用分离条件可以与恒等尺度的任务复杂度同时出现。
\end{counterexample}

这些例子共同表明，正文判据需要任务保真、状态压缩与闭合控制的联合证书。
